\documentclass[11pt,a4paper]{article}
% Preamble for CyclicTactic thesis/paper draft.
% Targets a generic article class; switch to a thesis class later by
% replacing \documentclass in main.tex.

\usepackage[utf8]{inputenc}
\usepackage[T1]{fontenc}
\usepackage{lmodern}
\usepackage[margin=1in]{geometry}
\usepackage{microtype}

% Math
\usepackage{amsmath, amssymb, amsthm, mathtools}
% Future-use: \usepackage{stmaryrd} (extra math symbols)
% Future-use: \usepackage{mathpartir} (inference rules)

% Map common Unicode chars in body text to math-mode equivalents so
% references like `\code{σ}` work without manual escaping.
\usepackage{booktabs}     % \toprule, \midrule, \bottomrule
\usepackage{pifont}       % \ding{52} for check marks

\usepackage{newunicodechar}
\newunicodechar{σ}{\ensuremath{\sigma}}
\newunicodechar{✓}{\ding{52}}
\newunicodechar{✗}{\ding{56}}
\newunicodechar{∀}{\ensuremath{\forall}}
\newunicodechar{∃}{\ensuremath{\exists}}
\newunicodechar{∈}{\ensuremath{\in}}
\newunicodechar{⊆}{\ensuremath{\subseteq}}
\newunicodechar{≤}{\ensuremath{\leq}}
\newunicodechar{≥}{\ensuremath{\geq}}
\newunicodechar{→}{\ensuremath{\to}}
\newunicodechar{↦}{\ensuremath{\mapsto}}
\newunicodechar{←}{\ensuremath{\leftarrow}}
\newunicodechar{§}{\S}

% Cross-references
\usepackage[colorlinks=true, linkcolor=blue!50!black, citecolor=blue!50!black,
            urlcolor=blue!50!black]{hyperref}
\usepackage{cleveref}

% Code listings (Lean + tactic syntax)
\usepackage{listings}
\usepackage{xcolor}
\definecolor{leankw}{rgb}{0.13,0.29,0.53}
\definecolor{leanstr}{rgb}{0.16,0.50,0.16}
\definecolor{leancmnt}{rgb}{0.45,0.45,0.45}
\lstdefinelanguage{Lean}{
  morekeywords={def,theorem,lemma,example,by,match,with,let,fun,end,
                namespace,structure,inductive,where,partial,instance,
                section,variable,open,import,return,do,if,then,else,
                cyclic_thm,cyclic_mutual,thm,cyclic,cyc_cases,back,
                end_mutual,refine,exact,apply,have,show,intro,intros,
                cases,induction,assumption,trivial,simp,rw,sorry,
                termination_by,decreasing_by},
  sensitive=true,
  morecomment=[l]{--},
  morecomment=[s]{/-}{-/},
  morestring=[b]",
  keywordstyle=\color{leankw}\bfseries,
  commentstyle=\color{leancmnt}\itshape,
  stringstyle=\color{leanstr},
  basicstyle=\ttfamily\footnotesize,
  showstringspaces=false,
  breaklines=true,
  columns=fullflexible
}
\lstset{language=Lean, frame=single, framerule=0.3pt}

% Theorem environments
\newtheorem{theorem}{Theorem}[section]
\newtheorem{lemma}[theorem]{Lemma}
\newtheorem{proposition}[theorem]{Proposition}
\newtheorem{corollary}[theorem]{Corollary}
\newtheorem{conjecture}[theorem]{Conjecture}

\theoremstyle{definition}
\newtheorem{definition}[theorem]{Definition}
\newtheorem{example}[theorem]{Example}

\theoremstyle{remark}
\newtheorem{remark}[theorem]{Remark}

% Macros
\newcommand{\Nat}{\mathbb{N}}
\newcommand{\Prop}{\mathrm{Prop}}
\newcommand{\Sort}{\mathrm{Sort}}
\newcommand{\Type}{\mathrm{Type}}
\newcommand{\code}[1]{\texttt{#1}}
\newcommand{\file}[1]{\texttt{\small #1}}
\newcommand{\Later}{\triangleright}
\newcommand{\Lob}{\textsf{L\"ob}}
\newcommand{\SCT}{\textsc{sct}}
\newcommand{\RelAnc}{\mathrm{RelAnc}}
\newcommand{\Ineq}{\mathrm{Ineq}}
\newcommand{\Hyp}{\mathrm{Hyp}}
\newcommand{\cov}{\mathrm{cov}}
\newcommand{\ResetGO}{\mathrm{Reset}}
\newcommand{\StackGO}{\mathrm{Stack}}
\newcommand{\NameGO}{\mathrm{Name}}


\title{Cyclic Proofs as Lean Tactics:\\
  A Paper-Faithful Implementation of\\
  Grotenhuis-Otten Unravelling}
\author{[Author]}
\date{\today}

\begin{document}
\maketitle

\begin{abstract}
We present \emph{CyclicTactic}, a Lean 4 library that lets users write cyclic
proofs as ordinary Lean tactics and translates them into kernel-checked Lean
theorems via a paper-faithful implementation of the Grotenhuis-Otten
Theorem~6.1 unravelling algorithm. The system records cyclic structure into a
side-channel \texttt{ProofTree} during tactic execution, validates it via the
size-change principle (Lee, Jones, and Ben-Amram, 2001), augments it with
Grotenhuis-Otten §5--§6 per-node data (Stack/Name/Reset annotation,
\(\mathrm{RelAnc}\)/\(\mathrm{Ineq}\)/\(\mathrm{Hyp}\)/\(\mathrm{cov}\)), and
emits a structural-induction proof that elaborates as the user-facing
declaration. Mutual cyclic systems are handled by an analogous emitter that
produces \texttt{mutual def \dots end} blocks with cross-theorem
\(\sigma\)-applied calls. We compare against Cyclist on Brotherston's
first-order benchmark and against the older Sprenger-Dam structural
emitter (\texttt{Unravel.translate}) we deleted in favour of the §6-faithful
path. We close by identifying the next algorithmic gap---non-structural
recursion as in merge-sort---and outlining two extensions (well-founded
emission, Löb induction) currently under separate development.
\end{abstract}

\tableofcontents

\section{Introduction}
\label{sec:introduction}

\subsection{Cyclic proofs in proof assistants}

Cyclic proof systems offer a natural representation for reasoning about
inductive predicates: rather than packaging an explicit induction hypothesis
into every recursive case, the proof tree may contain \emph{back-edges} that
point to ancestor nodes, provided a global termination criterion---typically
the size-change principle of Lee, Jones, and Ben-Amram~\cite{lee2001size}---is
satisfied. Sprenger and Dam~\cite{sprenger2003structure} established the
equivalence between cyclic proofs and ordinary inductive proofs in the
\(\mu\)-calculus, formally justifying the soundness of cyclic-proof systems
for first-order predicates. Brotherston's PhD thesis~\cite{brotherston2006phd}
and the Cyclist prover~\cite{brotherston2012cyclist} demonstrated the
practical viability of cyclic proof search for sequent calculi with explicit
trace conditions.

Despite this theoretical foundation, cyclic proofs have not seen widespread
adoption in mainstream interactive proof assistants. The gap is partly
ergonomic: cyclic proofs are usually presented in a sequent-calculus surface
syntax with explicit bud-companion edges, which sits awkwardly inside the
tactic-mode workflows that dominate Lean, Coq, and Isabelle. It is also
partly methodological: a cyclic proof's soundness depends on a
non-local termination condition (size-change), and integrating this with the
proof assistant's kernel-level type-checking is a nontrivial engineering
question.

A separate strand of theoretical work has recently appeared in
Grotenhuis and Otten's \emph{Unravelling Abstract Cyclic Proofs into Proofs
by Induction}~\cite{grotenhuis2026unravelling}. They give an algorithm,
applicable to any sufficiently abstract cyclic system, that translates a
cyclic proof into an inductive proof by following the local annotation
data they call the \emph{Stack/Name/Reset} representation (§5) and the
\emph{relativised ancestor / inequality / hypothesis / cover} per-node data
(§6). Their Theorem~6.1 then justifies a direct structural-induction emission
based on this data. Crucially, their algorithm operates on abstract cyclic
representations, so the translation is independent of the specific
proof-system or term-language details.

\subsection{Contributions}

This thesis presents \emph{CyclicTactic}, a Lean~4 library that lets users
write cyclic proofs as ordinary Lean tactics and, by implementing the
Grotenhuis-Otten Theorem~6.1 unravelling, replaces the user's interactive
cyclic-style proof with a kernel-checked structural-induction theorem.

Concretely, the contributions are:

\begin{enumerate}
  \item \textbf{Side-channel tactic-mode integration.} Three primitive
    tactics---\code{cyclic <label>}, \code{cyc\_cases x with}, and
    \code{back <label> \{σ\}}---let the user write cyclic proofs inside
    Lean's standard tactic mode. Real Lean goals are visible at every
    cursor position; cyclic structure is recorded into a side-channel
    \code{ProofTree} value via source-position attribution
    (§\ref{sec:architecture}).

  \item \textbf{A paper-faithful implementation of Grotenhuis-Otten §5--§6.}
    We implement the Stack/Name/Reset annotation step of Definition~5.1, the
    Lemma~5.9 key-bounded fixpoint unfolding, the Theorem~5.2
    progressing-name verification, and the §6 per-node augmentation
    (\(\RelAnc\), \(\Ineq\), \(\Hyp\), \(\cov\), \(\sigma\)) needed to
    drive emission. To our knowledge this is the first implementation of
    the Grotenhuis-Otten algorithm in any proof assistant.

  \item \textbf{Theorem 6.1 emission with canonical-form replacement.}
    The §6-driven emitter (\code{Theorem6.Emit.translate}) produces a
    Lean tactic script. After the user's cyclic proof elaborates as a
    recursive \code{def} (giving real InfoView goals during writing), the
    environment is rolled back and the §6-emitted theorem is re-elaborated as
    the canonical user-facing declaration.

  \item \textbf{Mutual cyclic systems.} A \code{cyclic\_mutual} surface
    command, together with an event demultiplexer and a mutual emitter
    (\code{Theorem6.Emit.translateMutual}), handles cyclic proofs that span
    multiple mutually-recursive theorems with cross-companion back-edges.
    Cross-theorem buds emit as direct sibling-theorem calls, relying on
    Lean's mutual-recursion termination check for soundness.

  \item \textbf{Comparison with Cyclist on Brotherston's first-order
    benchmark.} On the cases that both systems handle, our emitter produces
    proofs that match the structure Cyclist would emit. On
    Brotherston's tests~10 and~11 (associativity and commutativity of
    \texttt{ADD}, where Cyclist fails because its first-order proof search
    cannot perform equational reasoning), we succeed by decoupling the
    cyclic structure (validated by the size-change check) from the per-arm
    reasoning (closed by Lean's tactic ecosystem).
\end{enumerate}

\subsection{Outline}

§\ref{sec:background} reviews cyclic proof theory, the size-change principle,
and the Grotenhuis-Otten papers. §\ref{sec:architecture} describes the
side-channel architecture and the three primitive tactics.
§\ref{sec:implementation} walks through the implementation of the §5
annotation and §6 augmentation+emission. §\ref{sec:evaluation} reports the
empirical comparison against Cyclist and presents worked examples.
§\ref{sec:limitations} identifies the boundary cases the current system does
not handle---most prominently merge-sort, whose recursion is not
structural---and outlines two extensions currently under separate
development: well-founded emission via Lee 2009-style measure
synthesis~\cite{lee2009ranking}, and Löb-induction emission via a
step-indexed embedding~\cite{atkey2013productive}.

\section{Background}
\label{sec:background}

\subsection{Cyclic proofs}

Informally, a \emph{cyclic proof} of a sequent \(\Gamma \vdash \varphi\) is a
finite tree of inference steps in which each open branch (called a \emph{bud})
is decorated with a back-edge to an ancestor (called a \emph{companion})
whose sequent matches the bud's, modulo a substitution
\(\sigma\)~\cite{brotherston2006phd}. Cyclic proofs may be unsound in
general---one can construct trivial cyclic ``proofs'' of false
sequents---so a global termination criterion is required.

The standard criterion is Brotherston's \emph{trace condition}: every
infinite path through the cyclic proof's underlying graph must contain
infinitely many \emph{progressing} occurrences, where progression is defined
by a fixed structural-subterm relation. Lee, Jones, and Ben-Amram's
\emph{size-change principle} (SCT, \cite{lee2001size}) gives an algorithmic
incarnation of this idea: each back-edge is annotated with a
\emph{size-change graph} (SCG) recording which argument positions are
strictly smaller and which are non-strictly smaller; the system passes SCT iff
the closure of all SCG compositions contains, in every idempotent product, at
least one strict self-loop.

Sprenger and Dam's Theorem~5~\cite{sprenger2003structure} shows that any
cyclic proof satisfying the trace condition can be unfolded into a
tree-shaped proof using nested structural induction. This is the theoretical
basis for the \emph{structural} translation strategy that all
mainstream cyclic-to-inductive emitters (including Cyclist and our system)
implement.

\subsection{Cyclist and the first-order benchmark}

Cyclist~\cite{brotherston2012cyclist} is a standalone tool for cyclic
sequent-calculus proof search and first-order inductive reasoning. It
implements the trace condition directly and produces a separate
non-interactive proof artifact. Its \texttt{benchmarks/fo/} suite contains
tests~01--15, covering parity reasoning (\(\mathit{Od}\)/\(\mathit{Ev}\)
inductive predicates), addition properties, and various ordinals.

Cyclist's strengths are pure first-order proof search and explicit
trace-condition validation; its weaknesses are equational reasoning (it
cannot apply rewrite rules and so fails on the algebraic content of tests~10
and~11) and integration with mainstream proof
assistants. Our system inverts this trade-off: we lose Cyclist's automatic
proof search but gain Lean's tactic ecosystem for closing leaves and
intermediate steps.

\subsection{Grotenhuis-Otten unravelling}
\label{sec:bg-grotenhuis}

The Grotenhuis-Otten paper~\cite{grotenhuis2026unravelling} provides the
algorithmic core of our system. We summarise §5 and §6, which are the
sections we implement.

\paragraph{Abstract cyclic systems.}
The paper works in an abstract setting (§2): a cyclic proof is a tree whose
nodes are labelled with sequents and whose inference rules are taken from a
specified set, with each bud assigned a companion ancestor. The
\emph{induced inductive system} (§3) is the target of unravelling. The
authors give a general idea and case studies in §4 before introducing the
key constructions in §5--§6.

\paragraph{§5: Annotated cyclic representations.}
The paper defines (Definition~5.1) a per-node annotation consisting of:
\begin{itemize}
  \item a \(\StackGO\) function assigning each argument position a stack of
    fresh \emph{names},
  \item a \(\NameGO\) set tracking which names are currently alive, and
  \item a \(\ResetGO\) function recording which names have been ``reset'' by
    the most recent step.
\end{itemize}
The Definition~5.1 \emph{step} takes a parent annotation and an edge SCG and
produces the child annotation by (i) allocating fresh names per child slot,
(ii) reusing ancestor stacks where the SCG permits, (iii) running a reset
loop to fixpoint to capture the per-name dependency structure.

Lemma~5.4 bounds the alphabet at \(2^m + m\) names where \(m\) is the
arity, so by a Königs-lemma argument every cyclic proof has a finite
annotation reached after finitely many unfolding steps. Lemma~5.9 gives the
algorithmic shape of this unfolding: when a bud's initial annotation fails
to satisfy Theorem~5.2's progressing-name condition, iterate the annotation
up to a key-bounded fixpoint.

\paragraph{Theorem~5.2: Existence of a cyclic representation satisfying name
preservation.}
This theorem says the §5 annotation can be chosen so that every bud has a
\emph{progressing name}: a name in \(\ResetGO(\text{bud}) \cap \text{preserved
names along the cycle}\). Verifying Theorem~5.2 amounts to checking, for
each bud, that such a name exists.

\paragraph{§6: Relativised ancestor data and Theorem~6.1.}
Building on the §5 annotation, the paper defines per-node:
\begin{itemize}
  \item \(\RelAnc(n_k)\), the relativised-ancestor set;
  \item \(\Ineq(n_k)\), the inequalities the path enforces;
  \item \(\Hyp(n_k)\), the induction hypotheses in scope (one per
    \emph{sprout} ancestor that introduced a fresh IH);
  \item \(\cov(b)\) for each bud, the variable that ``covers'' the
    progressing name (i.e., binds the IH the bud applies);
  \item \(\sigma\) (Lemma~6.4), the substitution mapping sprout slot
    variables to bud slot variables and the progressing slot variable to
    \(\cov\).
\end{itemize}

Theorem~6.1 then states that the cyclic proof unfolds to a structural
induction proof in which:
\begin{itemize}
  \item every \emph{sprout} (a node that introduces a fresh IH in \(\Hyp\))
    emits as \texttt{induction <var> generalizing <rest> with};
  \item every other internal node emits as the corresponding C-rule
    (\texttt{cases} for case-splits, etc.);
  \item every bud emits as an IH application with arguments determined by
    \(\sigma\).
\end{itemize}

The paper continues with applications (§7), a multi-sort extension (§8), and
appendices on cyclic Heyting arithmetic (App.~B) and multi-sorted inductive
systems (App.~C). We implement §5 and §6 in full but do not address §7--§8
or the appendices.

\subsection{Related work in proof assistants}

To our knowledge, no mainstream proof assistant has previously shipped a
paper-faithful implementation of cyclic-to-inductive unravelling. The Iris
framework~\cite{jung2015iris} provides cyclic-style reasoning via the
\(\Later\) modality and Löb induction, but it sits at a different point in
the design space: Iris proofs live in a step-indexed embedding rather than
being translated into ordinary inductive proofs. The connection between
Iris-style Löb reasoning and cyclic proofs is a known
correspondence~\cite{simpson2017cyclic, berardi2017intuitionistic} and forms
the basis for the Löb-induction extension under separate development
(§\ref{sec:limitations}).

Wehr's PhD thesis~\cite{wehr2025cyclic} surveys the cyclic proof theory
landscape comprehensively and gives a fresh treatment of SCT-derived
induction orders in §3.2 (Theorem~3.2.4), which our system uses as an
auxiliary check alongside the Grotenhuis-Otten annotation.

The older Sprenger-Dam translation tradition is implemented in Cyclist and
in our pre-§6 baseline (\code{Unravel.translate}, since deleted). The
present work is the first to embed the Grotenhuis-Otten §6 algorithm itself
in a proof assistant.

\section{System architecture}
\label{sec:architecture}

This section describes the surface tactics, the side-channel architecture,
and the two-phase elaboration flow that lets users write cyclic proofs
interactively while producing a kernel-checked structural-induction theorem
as the user-facing declaration.

\subsection{Surface tactics}

CyclicTactic provides three primitive tactics, plus two command-level
forms.

\paragraph{\code{cyclic <label>}.}
Registers the current goal as a named \emph{companion}. Subsequent
back-edges within the same theorem body refer to this companion by label.
Internally, \code{cyclic} pushes a \code{.companion lbl seq} event into the
recorder (described below) and, for the first \code{cyclic} call in a body,
captures the goal head constant for use as \code{defaultSimpPred} when the
emitter needs to unfold the goal predicate.

\paragraph{\code{cyc\_cases x with | <pat> => \dots}.}
Case-splits on \code{x} while recording arm boundaries. Each arm's
syntactic range is captured so the recorder can later attribute back-edges
to the correct arm. The tactic delegates the actual case-split to Lean's
built-in \code{cases}, so InfoView goals appear as usual; the only
additional work is the event push.

\paragraph{\code{back <label> [\{σ\}]}.}
Closes the current goal via a recursive call to the enclosing theorem. The
\code{σ} substitution specifies how each binder of the enclosing theorem
maps to the call's arguments. \code{back R \{n := n', m := m\}} desugars to
\code{exact <name> n' m}, which Lean's standard structural-recursion check
accepts when \code{σ} produces strictly smaller arguments. The
back-edge's source position, σ, and surrounding source text are pushed as a
\code{.back} event.

\paragraph{Command-level forms.}
\code{cyclic\_thm name (binders) : type by tactics} packages the three
primitives into a command that produces a single user-facing declaration:
the kernel-checked structural-induction theorem produced by §6 emission, or
the fallback recursive form if §6 emission fails. \code{cyclic\_mutual
\dots end\_mutual} extends the same flow to mutually-recursive blocks
(§\ref{sec:impl-mutual}).

\subsection{Side-channel recording}

The cyclic structure of a proof---which case-splits, which back-edges, which
σs---is not visible to Lean's standard tactic machinery. Rather than
modifying Lean to track this information, CyclicTactic uses a
\emph{side-channel} architecture: each of the three primitive tactics pushes
an event into a global \code{CyclicState} \code{IO.Ref}; events are then
post-processed into a \code{ProofTree} value (a strongly-typed cyclic-proof
representation defined in \code{CyclicTactic/ProofTree.lean}).

The event-recording approach has three benefits over modifying Lean:
(1) the recorder is independent of Lean's elaborator internals, so it
survives Lean version upgrades; (2) the user gets ordinary InfoView goals
during proof writing, with no special editor support needed; (3) the
recorded structure is a pure-data Lean term, which we can validate, analyse,
and transform offline.

\begin{lstlisting}
inductive CyclicEvent where
  | companion (label : String) (sequent : Sequent)
  | caseSplitStart (var : String) (sequent : Sequent) (arms : List ArmInfo)
  | caseSplitEnd
  | back (ancestor : String) (sequent : Sequent) (σ : Subst)
         (pos : Nat) (sourceText : String)
\end{lstlisting}

The \code{ArmInfo} entries carry syntactic positions for each arm. The
\code{Build.eventsToTree} function (in \code{CyclicTactic/Build.lean})
walks the linear event stream: each \code{caseSplitStart}/\code{caseSplitEnd}
scope is converted to a \code{ProofTree.caseSplit}; back-edges within an arm
are mapped to that arm by source-position attribution.

\subsection{Two-phase elaboration}

The most subtle architectural choice is that \code{cyclic\_thm} elaborates
the proof \emph{twice} in succession: first as a recursive
\code{def}---which gives InfoView goals during interactive
writing---then, if §6 emission succeeds, as a structural-induction
\code{theorem} that becomes the user-facing declaration. The two phases
never coexist in the environment.

\begin{lstlisting}
-- Phase A: snapshot env before any work
let envBefore ← Lean.getEnv

-- Phase B: elaborate recursive form. User's tactics fire here;
-- InfoView shows real Lean goals at every cursor position.
Lean.Elab.Command.elabCommand
  (← `(def $name $binders : $type := by $tacs))

-- Phase C: snapshot env-with-recursive as fallback
let envWithRecursive ← Lean.getEnv

-- Phase D: build ProofTree from the recorded events;
-- run SCT + induction-order + §5 + §6 validations.
...

-- Phase E: try canonical §6 emission
let script := Theorem6.Emit.translate ...
match parse(script) with
| .ok cmdStx =>
  Lean.setEnv envBefore       -- roll back
  try elabCommand cmdStx       -- canonical form
  if hasGoodValue then ()      -- §6 emission committed
  else Lean.setEnv envWithRecursive  -- restore fallback
| .error _ =>
  -- recursive form stays as the user-facing decl
\end{lstlisting}

The user-facing \code{<name>} is always exactly one declaration. When §6
emission succeeds it is the structural-induction theorem; when §6 emission
fails for any reason it is the recursive \code{def} produced by Phase~B.

This architecture is what enables the system to be both ergonomic and
paper-faithful: the user writes cyclic-style with real interactivity, but
the kernel-checked artifact is the structural-induction proof Theorem~6.1
unravels. To our knowledge no prior cyclic-proof system in a mainstream
proof assistant uses this two-phase pattern.

\subsection{Pipeline overview}

The end-to-end flow from \code{cyclic\_thm} source to kernel-checked
theorem proceeds through nine stages, listed below with their owning
modules. \code{CyclicTactic/PIPELINE.md} walks the same flow on a
multi-recursive example (\code{btPredT}).

\begin{enumerate}
  \item \textbf{Command parsing and Phase~A env snapshot}
    (\code{Tactic.lean::elabCyclicThm}).
  \item \textbf{Phase~B recursive-form elaboration} (Lean's standard
    \code{def} machinery).
  \item \textbf{Event recording} (three custom tactics push into
    \code{CyclicState}).
  \item \textbf{ProofTree construction} (\code{Build.lean::eventsToTree},
    via source-position attribution).
  \item \textbf{SCT validation} (\code{ProofTree.extractTraceSCGsLabeled}
    \(+\) \code{SizeChange.checkMultiSCT}).
  \item \textbf{Paper-faithful §5 annotation}
    (\code{PaperAnnotation.lean}: Definition~5.1 step \(+\) Lemma~5.9
    unfolding \(+\) Theorem~5.2 verification).
  \item \textbf{§6 augmentation} (\code{Theorem6.lean::computeAug}:
    \(\RelAnc\)/\(\Ineq\)/\(\Hyp\)/\(\cov\)/\(\sigma\) per node).
  \item \textbf{§6 emission} (\code{Theorem6.Emit.translate}).
  \item \textbf{Canonical replacement} (Phase~E rollback and re-elaboration).
\end{enumerate}

\section{Implementation}
\label{sec:implementation}

This section describes the implementation of the §5 annotation, the §6
augmentation, the §6 emission for single-theorem proofs, and the extension
to mutual cyclic systems. Throughout we cite specific files and (where
relevant) function names in the repository.

\subsection{Definition 5.1 step and Lemma 5.9 unfolding}
\label{sec:impl-paper-annot}

The §5 annotation lives in \file{CyclicTactic/PaperAnnotation.lean}. We
implement:
\begin{itemize}
  \item \code{step : NodeAnnot \(\to\) SCGraph \(\to\) NodeAnnot}, the
    Definition~5.1 step.
  \item \code{freshNames : Nat \(\to\) List Name \(\to\) List Name} for the
    smallest-available naming convention.
  \item \code{resetLoop} for the reset-to-fixpoint operation.
  \item \code{annotateTree : ProofTree \(\to\) List BudAnnot}, which walks
    a proof tree and produces the per-bud annotation.
  \item \code{unfoldFixpoint}, which iterates \code{annotateTree} on a
    single failing bud up to a key-bounded fuel (16 iterations by
    default, well above the Lemma~5.4 bound for the proof sizes we have
    seen).
  \item \code{checkTheorem52}, which verifies the
    progressing-name condition.
\end{itemize}

The implementation closely follows the paper's prose. The two
non-obvious details are the smallest-available naming convention
(needed for the Lemma~5.4 alphabet bound to hold structurally) and the
key-based history tracking inside \code{unfoldFixpoint} (so the fixpoint
is detected when a bud's annotation \emph{key} recurs, not just its value).

\subsection{Per-node §6 augmentation}
\label{sec:impl-aug}

The §6 per-node data lives in \file{CyclicTactic/Theorem6.lean}. The
top-level entry point is:

\begin{lstlisting}
def computeAug (t : ProofTree) : NodeAugMap
\end{lstlisting}

returning, for each node label \(\ell\), an \code{NodeAug} record
containing \(\RelAnc\), \(\Ineq\), \(\Hyp\), and (for bud nodes) the
\code{bud?} info: the bud label, its progressing name, its cover, and the
progressing slot index.

The walker \code{computeAugAux} threads a current annotation
\code{cur~:~NodeAnnot} and a path substitution \code{σp~:~Subst} down the
tree. At each step it computes the parent's effective slots (the
\code{rootArgs} under \code{σp}), the child's effective slots (the
\code{rootArgs} under the child's \code{σc}, except at back-edges where the
child's literal sequent is used), and the edge SCG (a strict comparison
between the two slot lists, slot-by-slot).

The asymmetric treatment of back-edges versus other children is
load-bearing: at a back-edge, the bud's sequent already reflects the
user's \(\sigma\) substitution from the recursive call, so comparing it
against the parent's path-narrowed slots reveals strict descent. This
asymmetry, encoded in \code{effectiveSlots}, is what lets the §6 augmentation
detect the recursion structure the user wrote.

\subsection{§6 emission}
\label{sec:impl-emission}

The emitter \code{Theorem6.Emit.translate} walks the augmented tree and
produces a Lean tactic-script string. The structural choice per node is:

\begin{itemize}
  \item A \code{caseSplit} node \emph{at a sprout} (i.e., a node where
    \(\Hyp\) extends the parent's) emits as \texttt{induction <var>
    generalizing <rest> with}, where \texttt{<rest>} is the
    \code{rootArgs} except for the case-split variable.
  \item A \code{caseSplit} node \emph{not at a sprout} emits as
    \texttt{cases <var> with} (Lemma 6.3 C-rule selection).
  \item A \code{back} node emits as \texttt{exact ih\_<recVar> <args>},
    where \texttt{<recVar>} is read from \(\sigma\) at the prog slot and
    \texttt{<args>} are the σ-substituted rest-of-slots.
  \item A \code{node "branch"} (multi-recursive arm) emits as
    \texttt{simp [<pred>]; refine \(\langle\)?\_, \dots\(\rangle\); ·
    child; \dots} so that each branch's bud emits its IH call with the
    correct \(\cov\)-determined IH name.
  \item Leaves emit the user-supplied closing tactic (if any) or a default
    \code{simp}.
\end{itemize}

The decision about \texttt{induction} versus \texttt{cases} is the most
important fidelity gain over the pre-§6 emitter we deleted
(\code{Unravel.translate}, which always emitted \texttt{induction}). In
proofs where an inner case-split's subtree contains no bud that uses its
fresh IH, the §6 emitter correctly emits \texttt{cases}, avoiding an unused
\texttt{ih\_} binder.

The top-level function:

\begin{lstlisting}
def translate (defaultSimpPred : Option String) (goalType thmName : String)
    (varSorts : List (String × SortInfo)) (t : ProofTree) : String
\end{lstlisting}

calls \code{computeAug t} to build the per-node augmentation, then walks
the tree via \code{emitTree}, then assembles the result into a
\code{theorem <name> <binders> : <goalType> := by <body>} string ready to
hand to \code{Lean.Parser.runParserCategory}.

\subsection{Sort metadata and dependent binders}

The emitter needs to know each binder's inductive type to emit the right
constructor headers (e.g., \texttt{succ n'} for \texttt{Nat}, \texttt{node l r}
for binary trees). This information is stored as \code{SortInfo} records
(\file{CyclicTactic/EmitCommon.lean}) and built by
\code{Build.buildSortInfo} from a Lean \code{Expr}.

For binders whose types are non-dependent, the elaboration is
straightforward (\code{Lean.Elab.Term.elabType} on the binder's type
syntax). For \emph{dependent} binders like \code{(h : Od n)} where
\code{n} is an earlier binder, the per-binder \code{elabType} call would
fail because \code{n} is not in scope. We resolve this by reading
binder types from the \emph{already-elaborated recursive def's signature}
(via \code{Lean.Meta.forallTelescope}), which has the dependent context
properly built. This handling is in \code{Tactic.lean} immediately after
the §6 augmentation is computed.

\subsection{Canonical-form replacement}

The Phase~E replacement step is implemented as:

\begin{lstlisting}
match Lean.Parser.runParserCategory envWithRecursive `command script with
| .ok cmdStx =>
  Lean.setEnv envBefore                       -- roll back
  let beforeMsgs := (← getThe ...).messages
  try Lean.Elab.Command.elabCommand cmdStx
  catch _ => pure ()
  let envAfter ← Lean.getEnv
  let hasGoodValue : Bool :=
    match envAfter.find? name.getId with
    | some di => match di.value? with
                 | some v => !v.hasSorry
                 | none   => true
    | none    => false
  if hasGoodValue then ()                     -- §6 form committed
  else Lean.setEnv envWithRecursive           -- restore fallback
\end{lstlisting}

We check \code{hasSorry} rather than checking whether elaboration threw
an exception, because Lean's elaborator may emit error \emph{diagnostics}
(turning the value into one containing \code{sorry}) without raising; we
treat any sorry-containing value as a §6 emission failure.

\subsection{Mutual cyclic systems}
\label{sec:impl-mutual}

The \code{cyclic\_mutual} command extends the same flow to mutually-recursive
blocks. The implementation surfaces are:

\paragraph{Event demultiplexing.}
Inside a \code{cyclic\_mutual} block, events from all entries are
interleaved in the single event stream of \code{CyclicState}. The function
\code{demuxMutualEvents} (in \code{Tactic.lean}) walks the chronological
stream, using the pre-registered \emph{companion-target table} (each
\code{cyclic <label>} declaration's label \(\mapsto\) owning theorem) to
attribute each event to a specific theorem. The result is a list of
per-entry event subsets.

\paragraph{Per-entry tree construction.}
For each entry's events, \code{Build.eventsToTree} produces a
per-entry \code{ProofTree}. Each tree's back-edges reference companion
labels that may belong to \emph{other} entries (cross-companion edges).

\paragraph{Mutual emitter.}
\code{Theorem6.Emit.translateMutual} emits a \texttt{mutual def \dots end}
block. Cross-theorem buds (whose \code{ancestor} resolves to a sibling
entry's companion via the companion-target table) emit as
\texttt{exact <sibling-thm-name> <σ-applied args>}, relying on Lean's
mutual-recursion termination check for soundness. Within-entry case-splits
are emitted as \texttt{cases}, since the cyclic structure in the mutual
case is provided by Lean's mutual-recursion check, not by within-tree
sprouts.

\paragraph{§6 augmentation skipped for mutual.}
The §6 augmentation (\code{computeAug}) is skipped for mutual blocks
because \code{annotateTree}'s Lemma 5.9 unfolding does not terminate when
a bud's ancestor is not in the tree's ancestor stack (the cross-theorem
case). Since cross-theorem buds emit as direct sibling calls---not IH
applications---the §6 augmentation is unused in this case anyway. We
return to this limitation in §\ref{sec:limitations}: \code{cyclic\_mutual}
blocks with within-entry cycles are not yet supported.

\subsection{Module layout}

The implementation totals approximately 5\,400 lines of Lean code across
the following modules:

\begin{center}
\begin{tabular}{lr}
  \toprule
  Module & LOC \\
  \midrule
  \file{SizeChange.lean} & 171 \\
  \file{ProofTree.lean} & 400 \\
  \file{Extract.lean} & 125 \\
  \file{Measure.lean} & 285 \\
  \file{Annotation.lean} & 288 \\
  \file{InductionOrder.lean} & 307 \\
  \file{Reorganize.lean} & 369 \\
  \file{PaperAnnotation.lean} & 811 \\
  \file{Theorem6.lean} & 755 \\
  \file{EmitCommon.lean} & 155 \\
  \file{Tactic.lean} & 897 \\
  \file{Build.lean} & 372 \\
  \file{Design.lean} & 448 \\
  \midrule
  Total & 5\,383 \\
  \bottomrule
\end{tabular}
\end{center}

\section{Evaluation}
\label{sec:evaluation}

We evaluate CyclicTactic on three axes: (1) the worked-example suite under
\code{CyclicTactic/Examples/}; (2) Brotherston's first-order benchmark
\texttt{benchmarks/fo/} from the Cyclist
distribution~\cite{brotherston2012cyclist}; and (3) qualitative comparison
with the older Sprenger-Dam-style structural emitter we deleted in favour
of the §6 path.

\subsection{Worked-example suite}

The repository contains 18 \code{cyclic\_thm} and \code{cyclic\_mutual}
declarations distributed across five example files. All 18 elaborate via the
canonical §6 path (the \texttt{[cyclic\_thm <name>] canonical form: Theorem
6.1 emission ✓} diagnostic fires for each). The breakdown is:

\begin{center}
\begin{tabular}{lrl}
  \toprule
  File & Count & Coverage \\
  \midrule
  \file{Smoke.lean} & 3 & basic Nat induction, single back-edge \\
  \file{Probe.lean} & 5 & multi-arg, branches, \(\exists\)-witness, \texttt{have} \\
  \file{Probe2.lean} & 5 & enums, lists, nested back-edge \\
  \file{drp.lean} & 4 & comparison with the Cyclic DSL baseline \\
  \file{MutualSmoke.lean} & 1 & mutual block (Od/Ev) \\
  \bottomrule
\end{tabular}
\end{center}

The examples exercise every shape the system handles: single-recursive case
splits (e.g., \code{zeroAddT}, \code{addCommT}), multi-recursive branches
(\code{btPredT}), nested case-splits with lex descent (\code{ackTotal\_cyc}),
existential witnesses (\code{existsEqT}), intermediate \code{have} steps
(\code{zeroAddHaveT}), and cross-companion mutual blocks
(\code{od\_is\_nlike\_cyc} / \code{ev\_is\_nlike\_cyc}).

\subsection{Cyclist first-order benchmark}

Cyclist's \texttt{benchmarks/fo/} suite contains 15 numbered tests. We
ported the cases that map cleanly to Lean: tests 01 (mutual parity), 07
(commutative addition), 09 (\(\mathit{ADD}\) successor), 10 (associativity),
and 11 (commutativity). The results are summarised in
Table~\ref{tbl:cyclist}.

\begin{table}[h]
  \centering
  \caption{Cyclist comparison on \texttt{benchmarks/fo/} tests we ported.}
  \label{tbl:cyclist}
  \begin{tabular}{rlcc}
    \toprule
    \# & Sequent (cyclist notation) & Cyclist & CyclicTactic \\
    \midrule
    01 & \(O(x) \vdash N(x)\) (mutual) & ✓ & ✓ (\code{cyclic\_mutual}) \\
    07 & \(N(x) \vdash \mathit{ADD}(x, 0, x)\) & ✓ & ✓ \\
    09 & \(\mathit{ADD}(x,y,z) \vdash \mathit{ADD}(x, sy, sz)\) & ✓ & ✓ \\
    10 & associativity of \(\mathit{ADD}\) & ✗ & ✓ \\
    11 & commutativity of \(\mathit{ADD}\) & ✗ & ✓ \\
    \bottomrule
  \end{tabular}
\end{table}

Tests 10 and 11 are the headline win. Cyclist fails because its first-order
proof search applies sequent rules without equational simplification, so it
cannot manipulate intermediate \texttt{ADD}-predicate hypotheses
algebraically. Our system, hosted on Lean, decouples the cyclic structure
(validated by SCT and §6 augmentation) from the per-arm
reasoning (closed by Lean tactics), so \texttt{congrArg Nat.succ} plus
\texttt{rw} discharges the algebraic step inside each arm.

Tests 02, 04, 12, 13, 14, 15 are not ported in the current system. Tests
02 and 04 require introducing an induction hypothesis as a hypothesis and
case-splitting on its disjunctive result---our \code{back} primitive consumes
the goal directly via \code{exact}, so this pattern is not expressible
within the current tactic surface. Tests 12, 13, 15 require well-founded
emission with an explicit measure (Sprenger-Dam two-step descent, R/R2
lex, 2-Hydra)---our current §6 emitter is structural-only.

\subsection{Comparison with the deleted Unravel emitter}

The pre-§6 baseline (\code{Unravel.translate}, deleted as part of the §6
work) was a Sprenger-Dam-flavoured structural emitter that always emitted
\texttt{induction var with} for every case-split. The §6 emitter, by
contrast, distinguishes sprouts (case-splits whose subtree contains a bud
that uses the freshly-introduced IH) from non-sprouts (case-splits whose
subtree's buds use IHs from an enclosing sprout). At a non-sprout it
emits \texttt{cases var with} instead, dropping the unused
\texttt{ih\_} binder.

For 12 of the 17 single \code{cyclic\_thm} examples, both emitters produced
character-identical scripts. For the remaining 5 (\code{isColorT},
\code{fZeroT}, \code{maxNonneg}, \code{ackTotal\_cyc} inner split,
\code{existsEqT}, and \code{ackTotalT}), the §6 emitter chose \texttt{cases}
where the older emitter chose \texttt{induction}. In each case the §6 choice
is the paper-faithful one (Lemma 6.3's rule selection) and produces an
equally-elaborable proof with one fewer unused binder.

\subsection{Code-size and elaboration overhead}

The §6 emission adds approximately 1.5\,seconds of elaboration time per
\code{cyclic\_thm} on a modern laptop (this is dominated by the second
\code{Lean.Elab.Command.elabCommand} invocation in Phase~E). Total build
time of the worked-example suite (excluding the
\code{CyclistComparison.lean} file, which hangs on a pre-existing
\code{SubjectTerm.toString} stack overflow unrelated to this work) is under
30 seconds for all 18 examples combined.

The emitted proof terms are typically within a factor of 1.5 of the
hand-written structural-induction proof's term size, as measured by Lean's
\code{\#print} output. We have not yet measured this systematically across
the example suite.

\section{Limitations and ongoing extensions}
\label{sec:limitations}

\subsection{What the current system cannot express}

The current §6-structural emitter handles cyclic proofs whose recursion
fits Lean's structural-recursion check. Several classes of proofs lie
outside this scope.

\paragraph{Non-structural recursion.}
The most visible limitation is exposed by merge-sort. The auxiliary lemma
\(\code{length (merge xs ys) = length xs + length ys}\) recurses on the
result of a Decidable comparison: when \(x \le y\), it descends on
\code{xs}; otherwise on \code{ys}. Neither argument is structurally smaller
across both cases, so Lean's structural-recursion check rejects the
underlying \code{def} produced by Phase~B, and the structural §6 emitter
cannot witness termination on any single variable. This is precisely the
case the dispatcher (§\ref{sec:dispatch}) detects and routes to the
well-founded backend, where a lexicographic measure
\((\code{xs.length}, \code{ys.length})\) takes over; merge-sort additionally
exposes a residual arithmetic-discharge gap (the \code{+1} step the hand
proof leaves to \code{omega}) that is independent of the dispatch.

\paragraph{Decidable case-splits.}
The original \code{cyc\_cases} tactic recorded only inductive-type case
splits, so a proof body containing \code{by\_cases h : P} with back-edges
in each branch was opaque to the event recorder. The \code{cyc\_by\_cases}
tactic (\file{CyclicTactic/Tactic.lean}) now closes this gap: it records a
Decidable split as a \code{.dCaseSplitStart} event the same way
\code{cyc\_cases} records inductive ones, so the back-edges within each
branch enter the size-change check and the proof can dispatch to WF
emission (§\ref{sec:dispatch}).

\paragraph{DAG-shaped cyclic proofs.}
The system assumes back-edges target ancestors only. Cyclic proofs that
share intermediate sub-derivations across multiple branches (DAG-shaped
proofs in the sense of Brotherston's PhD~\cite{brotherston2006phd}) cannot
be expressed directly. The standard workaround---factoring shared
sub-derivations into separate \code{cyclic\_thm}s or \code{cyclic\_mutual}
entries---works for most practical proofs but is not always idiomatic. We
discuss the theoretical equivalence in §\ref{sec:dag-discussion} below.

\paragraph{Multi-companion in a single theorem.}
The §6 data model assumes companion = root: every back-edge in a single
\code{cyclic\_thm} body cycles to the same companion (the root sequent).
Proofs whose natural shape uses multiple \code{cyclic} declarations within
one theorem body (with different buds targeting different companions) are
not expressible. \code{cyclic\_mutual} offers a workaround by splitting
into mutually-recursive entries, but the surface ergonomics differ.

\paragraph{Within-entry cycles in \code{cyclic\_mutual}.}
The mutual emitter currently skips §6 augmentation, because
\code{annotateTree}'s Lemma 5.9 unfolding does not terminate on
cross-theorem buds. A \code{cyclic\_mutual} block whose entries contain
within-entry back-edges (in addition to cross-companion edges) is therefore
not supported. The §6 framework's extension to mutual systems is, to our
knowledge, an open theoretical question.

\paragraph{Proposition 5.8 verification.}
Our system uses Wehr~3.2.4's lex-order finder (Theorem~3.2.4 of
\cite{wehr2025cyclic}) as a check on the induction order, complementary to
the Grotenhuis-Otten Theorem~5.2 progressing-name verification. The paper's
Proposition~5.8 gives a stronger verification of the induction order that
we do not implement; we delegate this to Lean's structural-recursion check.
Implementing Proposition~5.8 directly would turn the current diagnostic
PASS/FAIL into a real proof obligation.

\paragraph{Multi-sort extension (§8).}
Grotenhuis-Otten §8 extends the unravelling algorithm to multi-sorted
inductive systems. We currently handle single-sorted cyclic proofs;
extending to §8 would let cyclic proofs span fundamentally distinct
inductive types as their recursion structure.

\subsection{Tree-shape versus DAG-shape: a discussion}
\label{sec:dag-discussion}

It is worth being precise about what is lost when we restrict to
tree-shaped cyclic proofs with back-edges to ancestors only. Sprenger and
Dam's Theorem~5 and its descendants establish that tree-shaped and
DAG-shaped cyclic proofs prove the same sequents: any DAG can be unfolded
to a tree by duplicating shared sub-derivations, with potentially
exponential blow-up. So provability is not affected.

What \emph{is} affected:

\begin{itemize}
  \item \textbf{Proof size.} A DAG with many shared sub-derivations
    unfolds to a tree exponentially larger.
  \item \textbf{Sharing-natural patterns.} Symmetry-driven proofs,
    cross-recursive lemmas with shared support, and confluence-style
    arguments naturally have DAG shape.
  \item \textbf{Lateral cycles.} A back-edge from one branch of a
    proof tree to a node in another branch (not an ancestor) cannot be
    represented.
\end{itemize}

For the Lean tactic-mode setting we target, this restriction is less
constraining than it sounds. Lean idiom favours factoring shared
sub-derivations into separate \code{theorem}s or local \code{have}
binders---both of which absorb the sharing-naturalness patterns above. The
remaining genuinely lateral case (a single proof tree with a non-ancestor
back-edge) is rare in practice and can typically be expressed by promoting
the lateral target to a separate \code{cyclic\_thm}.

This trade-off is most visible against Cyclist's separation-logic
benchmark (\texttt{benchmarks/sl/}), where DAG shape is heavily used by the
proof-search procedure and our system would struggle. For Cyclist's
first-order benchmark and for typical Lean proof work, tree-shape is a
defensible architectural choice rather than a deficiency.

\subsection{The dispatcher: structural, well-founded, and Löb emission}
\label{sec:dispatch}

The Grotenhuis-Otten §6 algorithm assumes structural emission is always
possible. It is not: merge-sort and the swap-recursion of Cyclist's
Test~14 (§\ref{sec:limitations}) admit no single-variable structural
order. The system therefore \emph{dispatches} among emission strategies on
the basis of the size-change data, rather than committing to structural
emission unconditionally. Two of the three strategies are now implemented
on the \code{wf} branch; the third remains under separate development. We
describe each, then state the characterisation theorem they converge on.

\paragraph{The dispatch predicate.}
The choice between structural and well-founded emission is made by a
single Boolean predicate, \code{canStructural}
(\file{CyclicTactic/Tactic.lean}):

\begin{lstlisting}
def canStructural (tree : ProofTree)
    (labeled : List (String × SCGraph)) (arity : Nat) : Bool :=
  match InductionOrder.findInductionOrder tree labeled arity with
  | none     => false
  | some asg => (InductionOrder.lexOrderFromAssignment asg).length <= 1
\end{lstlisting}

\code{findInductionOrder} runs the Wehr~3.2.4 lex-order finder over the
bud-companion strongly-connected components; it returns an assignment of
progressing argument positions, or \code{none} when no consistent order
exists. \code{canStructural} returns \code{true} exactly when that
assignment, after deduplication, uses \emph{at most one} argument
position---i.e.\ when a single \code{induction} on one variable carries the
whole recursion, which is precisely the shape Theorem~6.1's structural
emitter (§\ref{sec:impl-emission}) produces. When the order needs two or
more positions (a genuine lexicographic descent), or does not exist at all,
the predicate returns \code{false} and the proof is routed to the
well-founded backend. Merge-sort's \code{merge} fails the predicate
because its \(x\le y\) arm descends on \code{xs} (position~0) and its
\(x>y\) arm on \code{ys} (position~1): no single position is
preserved-and-progressing across both arms, so \code{findInductionOrder}
returns \code{none}.

\paragraph{Well-founded (WF) emission.}
When \code{canStructural} is \code{false}, the system emits
\code{def \dots\ termination\_by <measure>} with a
\code{decreasing\_by simp\_wf; omega} discharge, via
\code{WFEmit.translate} (\file{CyclicTactic/WFEmit.lean}). The measure is
synthesised from the per-back-edge size-change graphs by
\code{synthMeasure} (\file{CyclicTactic/Measure.lean}), which tries four
schemas in increasing generality (the lex-ranking ideas are
Lee~2009's~\cite{lee2009ranking}):

\begin{enumerate}
  \item \textbf{Lex over a permutation} of all argument positions---a
    strict self-loop at some position, nonstrict at all earlier ones
    (classic lex termination, e.g.\ Ackermann).
  \item \textbf{Lex over an ordered subset}, allowing non-participating
    positions to be dropped from the tuple.
  \item \textbf{Sum measure} \(a_0+\dots+a_{n-1}\), for swap-style
    recursions where the bag of arguments shrinks as a whole.
  \item \textbf{Closure-witness greedy lex}: compute the SCT composition
    closure, take each idempotent's strict-self-loop positions as a
    soundness witness, and greedily build a lex order covering the most
    uncovered idempotents at each step. This is a flat-arity
    specialisation of the Grotenhuis-Otten Definition~5.1 reset-proof
    construction.
\end{enumerate}

Every schema validates its candidate against the \emph{input} graphs, not
merely the closure idempotents, because Lean's \code{termination\_by}
requires every recursive call to decrease, not just every cycle. The
result is a \code{Measure}, which is either \code{lex order} or
\code{sum arity}. The Decidable case-splits that such proofs require
(merge-sort's \code{x \(\le\) y}) are now recorded by the \code{cyc\_by\_cases}
tactic, which emits a \code{.dCaseSplitStart} event exactly as
\code{cyc\_cases} does for inductive splits; this closes the gap noted
in §\ref{sec:limitations} above. The swap recursion of Test~14
(\code{N2(x,y)} with recursive call \code{N2(y,x)}) goes through
end-to-end: \code{synthMeasure} produces \code{sum 2}, rendered
\code{termination\_by x + y}, and the emitted \code{n2\_holds} is a
kernel-checked term (\file{CyclicTactic/Examples/WFSuite.lean}).

\paragraph{Löb-induction emission.}
A theoretically stronger alternative replaces structural and WF emission
entirely with \emph{Löb induction} via a step-indexed
embedding~\cite{atkey2013productive, jung2015iris}. Every SCT-validated
cyclic proof can be emitted as a Löb application, using a \(\Later\)
modality defined by \(\Later P\, n := \forall m < n, P\, m\) and the Löb
axiom \((\forall n.\, \Later P\, n \to P\, n) \to \forall n.\, P\, n\)
proved from \code{Nat.strongRec}. The advantage is universality: Löb
emission applies to any SCT-valid cyclic proof, with no case analysis on
which strategy fits. The disadvantage is artefact form: emitted proofs are
about \(\Later P\) rather than \(P\) directly, which is less idiomatic for
Lean users. This extension is under active development on a separate
\code{lob} branch in a sibling repository, and is the only one of the
three not yet integrated into the main dispatcher.

\paragraph{The characterisation theorem.}
The dispatcher \code{canStructural} and the synthesiser \code{synthMeasure}
are total, decidable functions on the §6 size-change data; what is
\emph{conjectural} is that their decisions agree with Lean's
own acceptance of the emitted artefact. We separate the two.

\begin{conjecture}[Dispatch soundness]
\label{conj:dispatch}
Let \(T\) be an SCT-valid \code{ProofTree} with labelled size-change
graphs \(G\) and arity \(n\), and write
\(A=\code{computeAug}\,T\) for its §6 augmentation.
\begin{enumerate}
  \item If \(\code{canStructural}\,T\,G\,n = \mathtt{true}\), then the
    structural emitter \(\code{Emit.translate}\,T\) produces a term Lean's
    structural-recursion check accepts.
  \item If \(\code{canStructural}\,T\,G\,n = \mathtt{false}\) and
    \(\code{synthMeasure}\,G\,n = \code{some}\,M\), then the well-founded
    emitter \(\code{WFEmit.translate}\,M\,T\) produces a \code{def} whose
    \(\code{termination\_by}\,M\) obligation \code{decreasing\_by} discharges.
  \item The residual case
    \(\code{canStructural} = \mathtt{false}\) and
    \(\code{synthMeasure} = \code{none}\)---an SCT-valid proof whose
    descent is captured by no lex-or-sum measure---is exactly the regime
    in which the first two strategies fail and Löb emission is required.
\end{enumerate}
\end{conjecture}

\begin{conjecture}[Löb universality]
\label{conj:lob}
For every SCT-valid \code{ProofTree} \(T\), the Löb emitter
\(\code{Emit.translateLob}\,T\) produces a Lean-acceptable term.
\end{conjecture}

The two extensions thus carve up the SCT-valid proofs into a structural
core (Conjecture~\ref{conj:dispatch}.1), a well-founded middle
(Conjecture~\ref{conj:dispatch}.2), and a Löb-only remainder
(Conjectures~\ref{conj:dispatch}.3 and~\ref{conj:lob}). Clauses~1 and~2
are now backed by working emitters and an example suite; proving that the
\emph{syntactic} predicates \code{canStructural} and \code{synthMeasure}
soundly approximate Lean's \emph{semantic} acceptance---and that the Löb
remainder is genuinely universal---are the load-bearing claims that anchor
the next stage of the work. A proof of Conjecture~\ref{conj:dispatch}.1
would proceed by showing the Wehr~3.2.4 single-position assignment is
exactly the witness Lean's structural-recursion elaborator searches for;
of clause~2, by appeal to the soundness of the Lee~2009 ranking schemas
against the input graphs; and of Conjecture~\ref{conj:lob}, by exhibiting
the step-indexed model in which \(\Later\) is the strict-order
approximation and reading SCT validity as the descent that makes the Löb
premise discharge.

\subsection{Other engineering follow-ups}

Several smaller items are not algorithmically interesting but are
worth listing:

\begin{itemize}
  \item The \code{CyclicTactic/Examples/CyclistComparison.lean} file
    has a pre-existing stack overflow in
    \code{SubjectTerm.toString}, triggered by one of the later examples
    in the file. The crash predates the §6 work and likely involves
    infinite term traversal when an inductive-predicate hypothesis is
    recorded. Investigating this is independent of the §6 / WF / Löb
    work but blocks bare \code{lake build}.

  \item The system does not yet support \code{cyclic\_thm} bodies whose
    structure is \emph{not} a sequence of \code{cyc\_cases} \(+\)
    \code{back} \(+\) leaves. In particular, intermediate \code{cyclic}
    re-declarations within a single body are not handled, and
    interleaving \code{cyclic\_thm}-style tactics with other custom
    tactics from libraries is untested.

  \item Diagnostic output is verbose but ad hoc; a structured
    \code{[cyclic\_thm]} info schema (so external tooling can parse the
    dispatcher's decisions) would make integration with editor extensions
    easier.

  \item Performance of the §6 augmentation has not been profiled. For
    very deep proof trees the \code{annotateTree} \(+\)
    \code{unfoldFixpoint} cost could become significant; we have not
    observed this on the example suite.
\end{itemize}


\bibliographystyle{plain}
\bibliography{bibliography}

\end{document}
